\documentclass[11pt,a4paper]{article}
\usepackage[margin=2.2cm]{geometry}
\usepackage{fontspec}
\setmainfont{TeX Gyre Pagella}
\usepackage{enumitem}
\usepackage{xcolor,listings}
\lstset{basicstyle=\ttfamily\small,frame=single,rulecolor=\color{gray!50},columns=fullflexible,keepspaces=true,breaklines=true}
\usepackage[hidelinks]{hyperref}
\setlist{nosep}
\setlength{\parskip}{4pt}

\title{Competitive Programming and Contests\\Hands-On 2 -- Report}
\author{Gabriele Marsili}
\date{November 2026}

\begin{document}\sloppy
\maketitle

The solutions are in \texttt{lib.rs} (\texttt{solve\_p1}, \texttt{solve\_p2} and the two segment trees),
together with the parsing of the input and the tests. \texttt{main.rs} only reads the input and prints the
answers, one per line. The package is named \texttt{handson2}, the name \texttt{main.rs} imports from:
\begin{lstlisting}
cargo run --release -- 1 < testsets/p1/input0.txt     # Problem 1
cargo run --release -- 2 testsets/p2/input0.txt       # Problem 2, input from a file
\end{lstlisting}
The input can be given on stdin, as the path of a file, or directly as command-line arguments. In the
following, $n$ is the size of the array (or the number of segments) and $m$ the number of queries.

\section*{Problem 1 -- Min and Max}

\textbf{Indices.} In the statement and in the test set the queries are 1-based, so \texttt{solve\_p1}
(line 226) converts them to 0-based before calling the tree.

\textbf{Data structure.} A segment tree (\texttt{SegmentTree}, line 112) stored in an array: the root is at
index 1 and the children of node $k$ are $2k$ and $2k+1$. Node $k$ covers a range $A[lo..hi]$ and stores the
maximum of that range. The array has $4n$ entries, which is enough for any $n$. The tree is built bottom-up
in $O(n)$.

\textbf{Max$(i,j)$.} The standard recursive visit: a node outside $[i,j]$ returns $-\infty$ (\texttt{i64::MIN}),
a node fully inside returns its value, and a node that partially overlaps asks both children and returns the
larger result. At most four nodes per level are visited, so the query costs $O(\log n)$.

\textbf{Update$(i,j,T)$.} Updating every position separately costs $O((j-i+1)\log n)$, i.e.\ $O(n \log n)$
per query in the worst case. To reach $O(\log n)$ I used lazy propagation. The key observation is that
taking the minimum with $T$ commutes with the maximum:
\[ \max_{p \in [lo,hi]} \min(A[p], T) = \min\Big(\max_{p \in [lo,hi]} A[p],\ T\Big). \]
So when the range of a node is fully inside $[i,j]$, its new value is simply $\min(t[k], T)$ and the
recursion can stop there. The children, which are not updated, are told later: \texttt{lazy[k]} stores the
pending value $T$ (\texttt{apply}, line 157). Two pending updates on the same node combine into a single one,
because $\min(\min(x, T_1), T_2) = \min(x, \min(T_1, T_2))$, so one value per node is enough.

Before a query or an update goes below a node, \texttt{push} (line 167) applies the pending value to the two
children and clears it. This is why \texttt{max} also takes \texttt{\&mut self}. A leaf is always either fully
inside or fully outside a range, so \texttt{push} is never called on a leaf. After the recursive calls, the
update recomputes the node as the maximum of its children (line 181).

\textbf{Observation:} if $t[k] \le T$, every value in the range is already at most $T$ and nothing changes,
so the update stops at that node without recording anything.

\textbf{Complexity:}
\begin{itemize}
\item $O(n)$ to build the tree, $O(\log n)$ per update and per query: $O((n+m)\log n)$ in total;
\item $O(n)$ space ($4n$ values and $4n$ pending updates).
\end{itemize}

\section*{Problem 2 -- Is There}

Positions and queries are 0-based. Let $\mathit{cover}[p]$ be the number of segments containing position $p$.
The query IsThere$(i,j,k)$ asks whether $\mathit{cover}[p] = k$ for some $p \in [i,j]$.

\textbf{Step 1: coverage.} The segments never change, so $\mathit{cover}$ is computed once with a difference
array (line 309): for each segment $\langle l, r\rangle$ I add $1$ at $l$ and $-1$ at $r+1$, and the prefix
sums give $\mathit{cover}[p]$. This costs $O(n)$. While computing it, I group the positions by their
coverage: \texttt{positions\_with[k]} is the list of positions $p$ with $\mathit{cover}[p] = k$ (line 318).
Since there are $n$ segments, $0 \le \mathit{cover}[p] \le n$.

\textbf{Step 2: queries offline, grouped by $k$.} A maximum or a sum over $\mathit{cover}$ does not answer
the question ``is the value $k$ present in the range'', so I answer the queries offline. I group the queries
by $k$ and use a segment tree over the positions (\texttt{CountTree}, line 241) that stores, for each node, how
many positions of its range are \emph{active}, with point update and range sum. Then, for each $k = 0,
\dots, n$ (line 330):
\begin{enumerate}
\item activate the positions in \texttt{positions\_with[k]};
\item for each query with this $k$, the answer is 1 if and only if the number of active positions in
$[i,j]$ is greater than 0;
\item deactivate the same positions, so the tree is empty again for the next $k$.
\end{enumerate}
Values of $k$ without queries are skipped. A query with $k > n$ is answered 0 directly, since no position is
covered by more than $n$ segments. The answers are stored by query index, so they are printed in the input
order.

\textbf{Complexity:}
\begin{itemize}
\item every position is activated and deactivated at most once ($O(n\log n)$), every query is one range sum
($O(m \log n)$), and grouping is $O(n+m)$: $O((n+m)\log n)$ in total;
\item $O(n+m)$ space: the coverage, the two groupings and the tree with $4n$ nodes.
\end{itemize}

\section*{Tests}

All the tests pass, the code is formatted with \texttt{cargo fmt} and \texttt{cargo clippy} gives no warnings.
\begin{itemize}
\item Official test set: a separate test runner (not part of the submission) compares the output of the binary with every
\texttt{outputK.txt}: 11/11 files for Problem 1 and 8/8 for Problem 2.
\item Unit tests in \texttt{lib.rs}: the examples of the statement for both problems; an update that covers
the whole array followed by a query on a single position (the pending value stored at the root must reach the
leaf); queries with $k > n$.
\item Randomized tests: both solutions are compared with brute-force versions on 3\,000 random inputs each
(small $n$, so that every branch of the recursion is exercised).
\item Parsing: the statement examples, input on a single line, truncated input, non-numeric tokens.
\item Performance: with $n = m = 200\,000$ the release build answers in 0.28 s (Problem 1, half of the
queries are updates on the whole array) and 0.13 s (Problem 2). The same input for Problem 1 took more than
60 seconds with one point update per position.
\end{itemize}

\section*{Sources}
\begin{itemize}
\item The statement of the Hands-On: \url{https://pages.di.unipi.it/rossano/blog/2024/handson2/}.
\item The lectures of the course on segment trees and lazy propagation.
\item CP-Algorithms, ``Segment Tree'', in particular the sections on range updates and lazy propagation:
\url{https://cp-algorithms.com/data_structures/segment_tree.html}.
\end{itemize}

\section*{Use of AI}
I used Claude Code to write the input parsing and the tests, to help me with lazy propagation and with the
solution of Problem 2, to review my code, and to translate into English and format this report.

\end{document}
